\documentclass{report}

\usepackage[spanish]{babel}
\usepackage[utf8x]{inputenc}
\usepackage{amsmath}
\usepackage{graphicx}
\usepackage[colorinlistoftodos]{todonotes}
\usepackage{enumitem}
\usepackage{listings}
\usepackage{verbatim}
\usepackage{eurosym}
\usepackage[export]{adjustbox}
\usepackage{amssymb}
\usepackage{bussproofs}
\usepackage{amsmath}
\usepackage{tikz}
\usepackage{xcolor}
\usepackage{listings}
\usepackage{titletoc}
\usepackage{hyperref}

\hypersetup{
  colorlinks=true,
  linkcolor=black,
  urlcolor=blue,
  citecolor=black
}

\newcommand{\coverPage}[6]{%
%----------------------------------------------------------------------------------------
%	COVER START
%----------------------------------------------------------------------------------------
\begin{titlepage}

    \newcommand{\HRule}{\rule{\linewidth}{0.5mm}}
    \newcommand{\department}{#1}
    \newcommand{\course}{#2}
    \newcommand{\titleValue}{#3}
    \newcommand{\subtitleValue}{#4}
    \newcommand{\authorName}{#5}

    \center

    %----------------------------------------------------------------------------------------
    %	HEADER
    %----------------------------------------------------------------------------------------
    \includegraphics{images/logo_usa.png}
    \vspace{0.5cm}
    \textsc{\Large \department}\\[0.5cm]
    \textsc{\Large \course}\\[0.5cm]
    \vfill

    %----------------------------------------------------------------------------------------
    %	TITLE
    %----------------------------------------------------------------------------------------

    \HRule\\
    \Huge
    \textbf{\titleValue}\\[0.5cm]
    \Large
    \textbf{\subtitleValue}\\
    \HRule\\[0.5cm]

    %----------------------------------------------------------------------------------------
    %	AUTHOR AND DATE
    %----------------------------------------------------------------------------------------

    \vfill
    \Large
    \textit{\authorName}\\
    {\large \today}\\[2cm]

\end{titlepage}
%----------------------------------------------------------------------------------------
%	COVER END
%----------------------------------------------------------------------------------------
}

\begin{document}
    \coverPage{ Matemáticas }{ Fundamentos de Matemáticas }{ Bono 3 }{   }{ Alexander Mendoza }{\today}
    \tableofcontents

    \pagebreak
    \chapter{ Bono 3 }

    \noindent Demuestre que no existen sucesiones estrictamente decrecientes de números naturales.\\

    Una sucesión $(x_n)$ es estrictamente decreciente si $x_n < x_m$ si $n > m$. Si la sucesión es finita, ciertamente existen sucesiones decrecientes en los números naturales, por ejemplo, \{ 5, 4, 3, 2, 1 \}. Sin embargo, si la sucesión es infinita contradice al principio del buen orden. La sucesión $(x_n)$ es un subconjunto de $\mathbb{N}$, por el principio del buen orden existe un $m \in (x_n)$ tal que $m < a$ para todo $a \in (x_n)$, lo cual es una contracción ya que $m$ sería el maximal de $(x_n)$, llamemos $x_a$ a $m$, luego $x_{a+1} < x_a$ lo cual es una contracción.\\

    \noindent Demuestre que los números racionales satisfacen la propiedad arquimediana. Propiedad Arquimediana: Si $x, y \in \mathbb{Q}$ y $x>0$, entonces existe un número entero positivo $n \in \mathbb{N}$ tal que $nx > y$.\\

    \noindent Sean $y$ y $x$ números racionales tal que $x > 0$, si $y \leq 0$ como $x > 0$ existe $n = 1 \in \mathbb{Q}$, así $x > yn$. Si $y > 0$, como $y$ y $x$ son racionales podemos reescribirlos como $y = \frac{a}{b}$ y $x = \frac{c}{d}$ donde $a, b, c \text{ y } d \in \mathbb{Z}$, luego existe $n \in \mathbb{Q}$ tal que $n > \frac{y}{x} = \frac{bc}{ad}$, así $nx > \frac{bc}{ad}x = \frac{bc}{ad} \cdot \frac{a}{b} = \frac{c}{d} = y$. De esta manera demostramos que los números racionales satisfacen la propiedad Arquimediana.
\end{document}
